\chapter{LA MATTINA E LA SERA}\label{la-mattina-e-la-sera}

\hfill\break

Guidai su per la collinetta che si affacciava sul ranch Condon e
parcheggiai in un punto in cui l'auto non poteva essere vista dalla
casa. Quando spensi i fari, riuscii a vedere il bagliore antelucano nel
cielo orientale, un'infausta schiaritura che portava via le nuove
stelle.

E fu allora che cominciai a tremare.

Non riuscivo a controllarmi. Aprii lo sportello e caddi fuori dall'auto
e mi risollevai grazie alla forza di volontà. Il terreno emergeva
dall'oscurità come un continente sommerso, brune colline, pascoli
trascurati restituiti al deserto, il lungo pendio che digradava
dolcemente verso la fattoria lontana. Mesquite e ocotillo tremavano nel
vento. Tremavo anch'io. Questa era paura: non la sparuta inquietudine
intellettuale con cui tutti avevamo imparato a convivere dall'inizio
dello Spin ma un panico viscerale, una paura simile a una malattia di
muscoli e budella. Tempo scaduto nel braccio della morte. Il giorno
della laurea. Carri e forconi in arrivo da oriente.

Mi chiedevo se Diane fosse spaventata come me. Mi chiedevo se sarei
riuscito a rincuorarla, se ancora mi rimanesse dentro un po' di
conforto.

Un'altra raffica di vento spazzò via la sabbia e la polvere dalla strada
arida del crinale. Forse il vento era il primo messaggero del Sole
rigonfio, un vento proveniente dalla faccia bollente del mondo.

Mi accovacciai in un punto da cui speravo di non essere visto e, ancora
tremante, riuscii a digitare il numero di Simon sul tastierino del
telefono.

Rispose dopo qualche squillo. Schiacciai il ricevitore sull'orecchio per
bloccare il rumore del vento.

«Non dovresti farlo» mi rimproverò.

«Sto interrompendo l'Estasi?»

«Non posso parlare.»

«Lei dov'è, Simon? In quale parte della casa?»

«\emph{Tu} dove sei?»

«Proprio sulla collina.» Il cielo diventava più luminoso, ogni secondo
di più, un livido porpora sull'orizzonte occidentale. Potevo vedere
chiaramente la fattoria. Non era cambiata un granché da quando c'ero
stato l'ultima volta, qualche anno prima. Il fienile esterno sembrava un
po' tirato a lucido, come se fosse stato riparato e imbiancato.

Ma la cosa molto più fastidiosa era la trincea scavata parallelamente al
fienile e coperta di tumuli di terra.

Un condotto fognario installato di recente, forse. Oppure una fossa
biologica. O una fossa comune.

«Sto venendo a trovarla» lo avvisai.

«Questo non è proprio possibile.»

«Suppongo che sia all'interno della casa. Una delle stanze da letto al
piano di sopra. Giusto?»

«Ma anche se la vedessi\ldots»

«Dille che sto arrivando, Simon.»

Sotto di me, vidi una sagoma che si muoveva tra la casa e il fienile.
Non era Simon. Non era Aaron Sorley, sempre che Fratello Aaron non
avesse perso una cinquantina di chili. Probabilmente era il Pastore Dan
Condon. Aveva un secchio d'acqua per mano. Sembrava avere fretta. Stava
succedendo qualcosa nel fienile.

«Stai rischiando la vita, qui» mi ammonì Simon.

Scoppiai a ridere. Era più forte di me.

Poi gli chiesi: «Sei nel fienile o in casa? Condon è nel fienile,
giusto? Che mi dici di Sorley e McIsaac? Come faccio a oltrepassarli?»

Sentii una pressione dietro il collo come una mano calda e mi girai.

Quella pressione era la luce solare. La corona del Sole aveva
attraversato l'orizzonte. La mia auto, le rocce, la fila disordinata di
ocotillo, tutte queste cose proiettavano una lunga ombra viola.

«Tyler? Tyler, non \emph{esiste} un passaggio percorribile. Devi\ldots»

Ma la voce di Simon fu soffocata da una scarica elettrostatica. La piena
luce solare doveva avere raggiunto l'aerostato che stava trasmettendo la
chiamata, tagliando il segnale. D'istinto, provai a richiamare ma il
telefono era inutilizzabile.

Rimasi accovacciato fino a quando il Sole salì per tre quarti,
guardandolo di striscio e poi distogliendo lo sguardo, ipnotizzato e
allo stesso tempo terrorizzato. Il disco era enorme e di un arancio
rossastro. Le macchie solari ci si arrampicavano come ferite infette. Di
tanto in tanto, veniva oscurato dagli sbuffi di polvere che si
sollevavano dal deserto circostante.

Poi mi alzai in piedi. Già morto, chissà. Magari fatalmente irradiato
senza nemmeno rendermene conto. Il caldo per il momento era
sopportabile, ma forse le mie cellule stavano subendo gravi danni a
causa dei raggi x che pungevano l'aria come proiettili invisibili.
Quindi rimasi in piedi e iniziai a camminare per la strada di terra
battuta verso la fattoria, in bella vista e senza armi. Proseguii
disarmato e indisturbato, almeno fino a quando giunsi nei pressi del
portico di legno. In altre parole, finché fratello Sorley, con tutti i
suoi centotrenta chili, si scaraventò fuori dalla porta zanzariera e mi
colpì sulla tempia con il calcio di un fucile.

ooo

Fratello Sorley non mi uccise, forse perché non voleva essere colto
dall'Estasi con le mani insanguinate. Invece mi buttò in una stanza
vuota al piano di sopra e chiuse la porta a chiave.

Passarono un paio d'ore prima di riuscire a tirarmi su senza conati di
vomito.

Quando si attenuarono le vertigini, andai alla finestra e alzai la
tendina di carta gialla. Da quell'angolo il Sole si trovava dietro la
casa, il terreno e il fienile erano lambiti da un violento bagliore
arancione. L'aria era già spietatamente calda ma, perlomeno, niente
stava andando a fuoco. Un gatto uscito dal fienile, noncurante
dell'esplosione nel cielo, lappava acqua stagnante da un fossato
all'ombra. Ipotizzai che il gatto sarebbe potuto campare fino al
tramonto. Forse anch'io.

Provai ad aprire la finestra - non che fossi proprio in grado di calarmi
da quell'altezza - ma era peggio che bloccata; le ante inchiodate, i
cardini saldati e il vecchio telaio evidentemente fermo in quella
posizione da quando era stato verniciato, anni prima.

Non c'erano mobili nella stanza a parte il letto, l'unico strumento che
avevo era quell'inutile telefono in tasca.

L'anta della porta era una lastra di legno massello e dubitavo di avere
la forza di abbatterla. Diane era probabilmente a pochi metri da me,
c'era solo una parete a dividerci. Ma non avevo modo di saperlo con
certezza né di provare a scoprirlo.

Anche il solo atto di pensare con lucidità alla situazione mi provocava
un dolore profondo e nauseante, localizzato dove il calcio del fucile mi
aveva fatto sanguinare la testa. Dovetti di nuovo sdraiarmi.

ooo

Prima della metà del pomeriggio il vento si era calmato. Quando
barcollando mi avvicinai di nuovo alla finestra riuscii a vedere il
bordo del disco solare che sovrastava la casa e il fienile, così grande
da sembrare perennemente in caduta, così vicino da poterlo quasi
toccare.

Rispetto al mattino la temperatura nella stanza al piano di sopra era
salita in fretta. Non avevo modo di misurarla ma immaginavo che fosse
intorno ai trentotto gradi buoni. Caldo ma non sufficiente a uccidere,
almeno non all'istante. Desiderai che ci fosse lì Jason a spiegarmela,
la termodinamica dell'estinzione globale. Forse mi avrebbe potuto
tracciare un grafico, mostrando dove le linee di tendenza convergevano
sulla mortalità.

Una foschia di calura vibrava verso l'alto dal suolo cotto dal Sole.

Dan Condon fece la spola su e giù dal fienile un altro paio di volte.
Era facile da riconoscere nella limpida intensità della luce diurna
arancione, aveva l'aria di uno del diciannovesimo secolo, la brutta
faccia squadrata e butterata con la barbetta: Lincoln con i blue jeans,
spilungone, determinato. Non alzò lo sguardo nemmeno quando colpii il
vetro.

Poi cominciai a tamburellare sulle pareti comunicanti pensando che Diane
avrebbe potuto picchiettare di risposta. Ma non ci fu nessuna replica.

Poi mi girò di nuovo la testa e caddi sul letto, l'aria nella stanza
chiusa era soffocante, il sudore inzuppava le lenzuola.

Mi addormentai, o persi conoscenza.

ooo

Mi svegliai pensando che la camera stesse andando a fuoco, ma era
soltanto la combinazione del caldo stagnante e di un tramonto fin troppo
acceso.

Tornai alla finestra.

Il Sole aveva attraversato l'orizzonte occidentale e stava colando a
picco in modo visibilmente rapido. Alte nuvole inconsistenti si
inarcavano nel cielo all'imbrunire, pezzi di umidità succhiati da un
terreno già prosciugato. Vidi che qualcuno aveva portato la mia auto giù
dalla collina e l'aveva parcheggiata proprio a sinistra del fienile. E
aveva preso le chiavi, senza dubbio. Non che contenesse abbastanza
benzina da poterla utilizzare.

Comunque ero sopravvissuto a quella giornata. Pensai: ``\emph{Siamo}
sopravvissuti''. Entrambi. Io e Diane. E di sicuro milioni di altri.
Quindi questa era la versione lenta della fine del mondo. Ci avrebbe
ucciso cucinandoci un grado alla volta; oppure, se non avesse
funzionato, radendo al suolo l'ecosistema terrestre.

Il Sole rigonfio finalmente sparì. L'aria sembrò subito più fresca di
una decina di gradi.

Qualche sporadica stella apparve tra le nuvole diafane.

Non avevo mangiato e avevo una sete terribile. Forse il piano di Condon
consisteva nel lasciarmi qui a morire disidratato\ldots{} o forse si era
semplicemente dimenticato di me. Non riuscivo nemmeno a immaginare come
il Pastore Dan stesse formulando tutto questo nella sua mente, se si
sentisse giustificato o terrorizzato o una qualche combinazione tra le
due cose.

La stanza si fece buia. Nessuna illuminazione dall'alto, nessuna
lampada. Ma potevo sentire un debole scoppiettio che doveva provenire da
un generatore di corrente a benzina, e la luce filtrava dalle finestre
del primo piano e dal fienile.

E invece io non possedevo niente di tecnologico tranne il telefono. Lo
tirai fuori dalla tasca e lo accesi, svogliato, solo per vedere la
luminescenza dello schermo.

Poi pensai un'altra cosa.

ooo

«Simon?»

Silenzio.

«Simon, sei tu? Riesci a sentirmi?»

Silenzio. Poi una voce metallica, digitalizzata: «Mi hai fatto quasi
morire di paura. Pensavo che quest'affare fosse rotto.»

«Solo di giorno, alla luce del Sole.»

L'interferenza solare aveva impedito le trasmissioni dagli aerostati
d'alta quota. Ma ora la Terra ci stava schermando dal Sole. Forse gli
aerostati avevano subito dei danni - il segnale suonava a banda bassa e
statica - ma per il momento il ritorno era abbastanza buono.

«Mi spiace per ciò che è successo,» si scusò, «ma ti avevo avvisato.»

«Dove sei? Casa o fienile?»

Pausa. «Casa.»

«Vi ho osservato per tutto il giorno e non ho visto la moglie di Condon
e nemmeno quella di Sorley con i figli. Neppure McIsaac o la sua
famiglia. Che fine hanno fatto?»

«Sono andati via.»

«Ne sei sicuro?».

«Se ne sono sicuro? Certo. Diane non è stata l'unica ad ammalarsi. Solo
l'ultima. La prima è stata la figlia piccola di Teddy McIsaac. Poi il
figlio, poi Teddy stesso. Quando ha capito che i suoi figli
erano\ldots{} be', capisci, evidentemente molto malati, e non
miglioravano, be', a quel punto li ha caricati sul furgone e se n'è
andato. La moglie del Pastore Dan li ha seguiti.»

«Quando è successo?»

«Un paio di mesi fa. La moglie di Aaron e i suoi figli sono partiti da
soli poco dopo. La loro fede vacillava. E poi avevano paura di
ammalarsi.»

«Li hai visti andare via? Ne sei certo?»

«Sì, perché non dovrei esserlo?»

«Nella fossa accanto al fienile sembra proprio che ci sia sepolto
qualcosa.»

«Ah, quella! Be', sì, hai ragione, \emph{c'è} qualcosa sepolto
laggiù\ldots{} il bestiame nocivo.»

«Scusa?»

«Un uomo di nome Boswell Geller possedeva un grande ranch sulla Sierra
Bonita. Era un amico del Jordan Tabernacle prima della riorganizzazione.
Ed era amico del Pastore Dan. Allevava giovenche rosse, ma alla fine
dell'anno scorso il Dipartimento dell'Agricoltura ha fatto partire
un'indagine. Proprio quando stava facendo progressi! Boswell e il
Pastore Dan volevano fare riprodurre tutte le varietà rosse di bestiame
del mondo, perché potessero rappresentare la conversione dei Gentili. Il
Pastore Dan dice che è proprio di questo che si parla in Numeri 19 - una
rossa giovenca pura nata alla fine del tempo, da razze provenienti da
ogni continente, ovunque sia stato predicato il Vangelo. Il sacrificio è
allo stesso tempo letterale e simbolico. Nel sacrificio biblico le
ceneri della giovenca hanno il potere di purificare una persona
corrotta. Ma alla fine del mondo il Sole consuma la giovenca e le ceneri
vengono sparse nei quattro punti cardinali, purificando l'intera Terra
dalla morte. Questo è ciò che sta succedendo ora. Ebrei 9:
``\emph{Perché, se il sangue dei capri e dei vitelli e la cenere di una
	giovenca sparsi su quelli che sono contaminati, li santifica,
	purificandoli nella carne, quanto più il sangue di Cristo che mediante
	lo Spirito eterno ha offerto se stesso puro d'ogni colpa a Dio,
	purificherà la nostra coscienza dalle opere morte per servire il Dio
	vivente}?'' Quindi, ovviamente\ldots»

«Avete conservato quel bestiame \emph{qui}?»

«Solo alcuni. Quindici animali da monta sono stati fatti sparire prima
che il Dipartimento dell'Agricoltura potesse rivendicarli.»

«Da quel momento le persone hanno iniziato ad ammalarsi?»

«Non solo le persone. Anche il bestiame. Abbiamo scavato quella fossa
accanto al fienile per seppellirli, tutti tranne tre esemplari del ceppo
originario.»

«Spossatezza, portamento instabile, perdita di peso prima della morte?»

«Sì, nella maggior parte dei casi\ldots{} come fai a saperlo?»

«Questi sono i sintomi della sindrome da deperimento cardiovascolare. Le
mucche erano portatrici. Ecco cosa non va in Diane.»

Seguì un lungo silenzio. Poi Simon disse: «Non posso parlare di questo
con te.»

Risposi: «Sono al piano superiore nella stanza sul retro\ldots»

«So dove ti trovi.»

«Allora vieni e apri la porta.»

«Non posso.»

«Perché? Sei sorvegliato da qualcuno?»

«No, è che non posso proprio liberarti. Non dovrei nemmeno parlare con
te. Ho da fare, Tyler. Sto preparando la cena per Diane.»

«Ha ancora la forza di mangiare?»

«Un po'\ldots{} se la aiuto.»

«Fammi uscire. Non devono saperlo per forza.»

«Non posso.»

«Ha bisogno di un medico.»

«Non potrei farti uscire neppure se volessi, fratello Aaron ha le
chiavi.»

Valutai la situazione. «Allora quando le porti la cena lasciale il
telefono - il tuo telefono. Hai detto che voleva parlare con me,
giusto?»

«La metà delle volte dice cose che non pensa.»

«Credi che fosse una di quelle?»

«Non posso più parlare.»

«Lasciale il telefono, Simon. Simon?»

Linea morta.

ooo

Andai alla finestra, osservai e aspettai.

Vidi il Pastore Dan portare dal fienile alla casa due secchi vuoti, poi
tornare indietro con i due secchi pieni e fumanti. Alcuni minuti dopo
Aaron Sorley lo raggiunse.

E quindi erano rimasti in casa solo Simon e Diane. Forse le stava
servendo la cena. Imboccandola.

Non vedevo l'ora di usare il telefono, ma mi ero ripromesso di
aspettare, lasciare che le cose si sistemassero un po', che la
situazione si calmasse.

Guardai il fienile. Una luce intensa filtrava tra le assi delle pareti
come se qualcuno avesse installato una rastrelliera di fari industriali.
Condon aveva fatto avanti e indietro per tutto il giorno. In quel
fienile stava succedendo qualcosa, ma Simon non mi aveva detto cosa.

La luce fioca del mio orologio segnalò che era passata un'ora.

Poi sentii un rumore indistinto, poteva essere una porta che si
chiudeva, passi sulle scale; e un attimo dopo vidi Simon dirigersi verso
il fienile.

Non guardò in alto.

E una volta entrato non ne uscì. Era all'interno con Sorley e Condon, e
se aveva ancora con sé il telefono, e se era stato tanto stupido da aver
attivato la suoneria, chiamarlo in quel momento lo avrebbe messo in
pericolo. Non che fossi particolarmente preoccupato del benessere di
Simon.

Ma se aveva lasciato il telefono a Diane quello era il momento.

Digitai il numero.

«Sì,» disse - era la voce di Diane - e poi l'inflessione salì di tono, e
divenne domanda: «Sì?»

La sua voce era ansimante e tenue. Quelle due sillabe erano già
sufficienti per richiedere una diagnosi.

«Diane. Sono io. Sono Tyler» mormorai cercando di tenere a freno le mie
stesse pulsazioni rabbiose. Sentivo come se mi si fosse spalancato un
varco nel petto.

«Tyler,» bisbigliò, «Ty\ldots{} Simon me l'ha detto, che forse avresti
chiamato.»

Dovetti sforzarmi per capire le sue parole. Non c'era nessuna forza a
spingerle; solo gola e lingua, niente torace. Il che era coerente con
l'eziologia della sindrome da deperimento cardiovascolare. La malattia
colpiva per primi i polmoni, poi il cuore, con un attacco simultaneo
dall'efficacia quasi militaresca. Il tessuto polmonare cicatrizzato e
schiumoso trasmetteva meno ossigeno al sangue; il cuore, in carenza di
ossigeno, pompava il sangue in modo meno efficiente; il batterio della
sindrome sfruttava entrambe le debolezze, azzannando il corpo a ogni
respiro faticoso.

«Non sono lontano» le dissi. «Sono davvero vicino, Diane.»

«Vicino. Posso vederti?»

Avrei voluto fare un buco nel muro. «Presto. Promesso. Dobbiamo portarti
via da qui. Soccorrerti. Rimetterti in sesto.»

Sentii che il suo respiro si faceva più agonizzante e mi chiesi se avevo
perso la sua attenzione. Poi sussurrò: «Penso di aver visto il
Sole\ldots»

«Non è la fine del mondo. Non ancora, comunque.»

«Non ancora?»

«No.»

«Simon» disse.

«Cosa c'entra lui?»

«Ne sarà così deluso.»

«Hai la sindrome da deperimento cardiovascolare, Diane. La stessa che
probabilmente aveva la famiglia di McIsaac. Hanno fatto bene a farsi
aiutare. È una malattia curabile.» Non aggiunsi ``fino a un certo
punto'' o ``sempre che non sia progredita fino allo stadio terminale''.
«Ma dobbiamo portarti via da qui.»

«Mi sei mancato.»

«Anche tu mi sei mancata. Capisci quello che sto dicendo?»

«Sì.»

«Sei pronta ad andare via?»

«Se arriverà il momento\ldots»

«Quel momento è molto vicino. Fino ad allora riposati. Ma può darsi che
ci dovremo dare una mossa. Lo capisci, Diane?»

«Simon» mormorò con un filo di voce. «Deluso» aggiunse.

«Tu riposati, e io\ldots»

Ma non ebbi il tempo di finire la frase.

Una chiave sferragliò nella porta. Chiusi con uno scatto il cellulare e
lo infilai in tasca. La porta si aprì e Aaron Sorley si piantò sulla
soglia, fucile in mano, ansimando come se avesse corso su per le scale.
Era una sagoma delineata dalla luce fioca che proveniva dal corridoio.

Arretrai fino a ritrovarmi con le spalle al muro.

«L'etichetta sulla tua patente afferma che sei un medico» ringhiò. «È
giusto?»

Annuii.

«Allora vieni con me.»

ooo

Sorley mi fece marciare giù dalle scale e fuori dalla porta posteriore
verso il fienile.

La Luna, velata dal colore ambrato della luce del Sole gibboso,
sfigurata e più piccola di quanto mi ricordassi, si era arrampicata
sull'orizzonte orientale. L'aria notturna era fresca, quasi inebriante.
Feci respiri profondi. Il sollievo durò fin quando Sorley aprì sbattendo
la porta del fienile e ne fuoriuscì un grezzo fetore animale - puzza di
mattatoio, sangue ed escrementi.

«Entra» grugnì Sorley, dandomi una spinta con la mano libera.

La luce proveniva da una grossa lampada ad alogenuri collegata ad un
cavo d'alimentazione e posta sopra uno stabbio. Un generatore a benzina,
nascosto da qualche parte, sferragliava emettendo un suono simile a
qualcuno che sgassa in lontananza con una motocicletta.

Dan Condon stava sul lato aperto dello stabbio, inzuppando le mani in un
secchio d'acqua bollente. Quando entrammo, alzò lo sguardo. Mise il
broncio, la sua faccia era una geografia desolata sotto lo sfolgorante
occhio-di-bue della luce, ma sembrava meno minaccioso di quanto
ricordassi. Infatti, pareva assottigliato, macilento, forse addirittura
malato, forse nella fase iniziale del suo personale caso di sindrome da
deperimento cardiovascolare. «Richiudi quella porta» ordinò.

Aaron la chiuse con una spinta. Simon, a pochi passi da Condon, mi
lanciava occhiate rapide e nervose.

«Vieni qui» mi disse Condon. «Abbiamo bisogno del tuo aiuto. Delle tue
competenze mediche, se possibile.»

Nello stabbio, su un lurido letto di paglia, una giovenca rinsecchita
stava cercando di partorire un vitello.

La giovenca giaceva a terra, la groppa ossuta sporgeva dal recinto. La
coda le era stata legata al collo con un pezzo di spago per evitare che
fosse d'intralcio. Il sacco amniotico sporgeva dalla vulva, la paglia
attorno era punteggiata di muco insanguinato.

Mi ritrassi: «Non sono un veterinario.»

«Lo so» ribatté Condon. Nei suoi occhi c'era un'isteria soffocata, lo
sguardo di un uomo che ha dato una festa di cui ha perso il controllo:
gli ospiti diventati selvaggi, i vicini che si lamentano, bottiglie di
liquore che volano dalle finestre come colpi di mortaio. «Ma ci serve
un'altra mano.»

Tutto ciò che sapevo sulla riproduzione e sul parto degli animali
l'avevo imparato dalle storie di Molly Seagram sulla vita nella fattoria
dei suoi genitori. Nessuna delle storie era stata particolarmente
gradevole. Perlomeno Condon aveva provveduto a tutte le basi necessarie,
per quel che ricordavo io: acqua calda, disinfettante, catene
ostetriche, una grossa bottiglia d'olio minerale già chiazzata con
impronte di mani insanguinate.

«È in parte Angler,» mi informò Condon, «in parte Danish Red, in parte
Belarus Red, e questa è solo la sua discendenza più recente. Ma
l'ibridazione è un fattore di rischio per la distocia, lo diceva sempre
fratello Geller. La parola ``distocia'' significa travaglio difficile. I
meticci hanno spesso problemi di parto. È in travaglio da almeno quattro
ore. Dobbiamo estrarre il feto.»

Condon espresse tutto ciò con tono distaccato, come un uomo che tiene
una lezione per una classe d'idioti. Non sembrava importante chi fossi
né come fossi arrivato lì, l'unica cosa rilevante era che avevano a
disposizione un aiuto in più.

«Ho bisogno d'acqua» dissi.

«Per lavarsi c'è un secchio.»

«Non dico per lavarmi. Non bevo niente da ieri notte.»

Condon fece una pausa come per valutare questa informazione. Poi annuì.
«Simon. Provvedi.»

Simon sembrava il ragazzo di bottega del trio. Chinò il capo ed esclamò:
«Ti porto io da bere, Tyler, tranquillo» cercando ancora di evitare di
guardarmi mentre Sorley apriva la porta del fienile per farlo uscire.

Condon si voltò di nuovo verso lo stabbio, dove la giovenca esausta
giaceva ansimante. Aveva i fianchi decorati da mosche indaffarate.
Alcune ronzavano ignorate sulle spalle di Condon. Si impregnò le mani di
olio minerale e si accucciò per espandere il canale del parto della
giovenca, la faccia contorta in un misto di fervore e disgusto. Ma aveva
appena iniziato quando il vitello venne annaffiato dall'ennesimo
zampillo di sangue e fluidi, la testa che emergeva appena nonostante le
impetuose contrazioni della giovenca. Era troppo grosso. Molly mi aveva
parlato di vitelli di dimensioni eccessive - non una situazione grave
come un parto podalico o un blocco dell'anca, ma comunque sgradevole da
affrontare.

Non era certo d'aiuto il fatto che la giovenca fosse chiaramente
ammalata, sbavasse muco verdastro e faticasse a respirare anche quando
le contrazioni si calmavano. Non sapevo se fosse il caso di parlarne a
Condon. Era evidente che anche la sua divina giovenca era stata
contagiata.

Ma il Pastore Dan non lo sapeva o non gli importava. Condon era l'unico
superstite dell'ala dispensazionalista del Jordan Tabernacle, una chiesa
dentro la chiesa, ridotta a due parrocchiani, Sorley e Simon, e potevo
solo immaginare quanto dovesse essere stata salda la sua fede per
sorreggerlo sul cammino verso la fine del mondo. Esclamò, con il solito
tono di isteria soffocata: «Il vitello, il vitello è rosso\ldots{}
Aaron, guarda il vitello.»

Aaron Sorley, che era appostato sulla porta col fucile, si avvicinò per
guardare nello stabbio. Il vitello era davvero rosso. Impregnato di
sangue. Era anche floscio.

Sorley chiese: «Respira?»

«Lo farà» rispose Condon. Era assorto, sembrava assaporare quel momento
in cui genuinamente credeva che il mondo stesse per incardinarsi
sull'eternità. «Mettete le catene attorno ai metacarpi, veloci, ora!»

Sorley mi lanciò un'occhiata che fu anche un avvertimento - del tipo:
``non ti azzardare a dire nemmeno una parola, cazzo'' - e seguimmo le
istruzioni, lavorando fino a quando fummo impiastricciati di sangue fino
ai gomiti.

L'atto di far nascere un vitello sovradimensionato è violento e assurdo,
un grottesco connubio di biologia e forza bruta. Per un parto del genere
occorrono almeno due uomini discretamente forti. Le catene ostetriche
servono a tirare e gli strappi devono essere a tempo con le contrazioni
della giovenca, altrimenti l'animale rischia di essere eviscerato.

Ma quella bestia era allo stremo delle forze, e il suo vitello - la cui
testa ciondolava senza vita - era a quel punto un feto morto.

Guardai Sorley e Sorley guardò me. Rimanemmo in silenzio. Quindi, Condon
disse: «La prima cosa che dobbiamo fare è tirarlo fuori. Poi lo
rianimiamo.»

Dalla porta del fienile soffiò una corrente d'aria più fresca. Era
Simon, di ritorno con una bottiglia d'acqua di sorgente. Fissò noi e poi
il feto morto partorito per metà, la faccia pallida e l'espressione
attonita.

«Ti ho preso da bere.» Fu tutto quello che riuscì a dire.

La giovenca ebbe un'altra contrazione fiacca e improduttiva. Lasciai
cadere la catena. Condon disse: «Tu bevi, figliolo. Poi riprendiamo.»

«Devo ripulirmi. Almeno lavarmi le mani.»

«L'acqua calda pulita è nei secchi vicino alle balle di fieno. Ma fai
alla svelta.» Aveva gli occhi serrati, stretti, per concentrarsi sulla
battaglia interiore tra buon senso e la sua fede.

Mi lavai e disinfettai le mani. Sorley mi osservava attentamente. Le sue
mani stringevano la catena ostetrica, ma il fucile era poggiato al
corrimano dello stabbio, a portata di mano.

Quando Simon mi porse la bottiglia, mi appoggiai alla sua spalla e gli
sussurrai: «Non posso aiutare Diane a meno di portarla via da qui. Lo
capisci? E non posso farlo senza il \emph{tuo} aiuto. Serve un mezzo di
trasporto affidabile, il serbatoio pieno, e Diane deve essere a bordo,
preferibilmente prima che Condon si renda conto che il vitello è morto.»

Simon rimase a bocca aperta. «È davvero morto?» chiese a volume troppo
alto, ma né Sorley né Condon sembrarono accorgersene.

«Il vitello non respira» risposi. «La giovenca è a malapena viva.»

«Ma il vitello è rosso? Tutto rosso? Niente macchie bianche o nere? Solo
rosso?»

«Anche se fosse un dannato camion dei pompieri, Simon, non aiuterebbe
Diane.»

Mi guardò come se gli avessi annunciato che il suo cucciolo era stato
investito. Mi chiesi quando avesse barattato la fiducia in se stesso con
questo vuoto disorientamento, se fosse successo all'improvviso e la
gioia si fosse prosciugata un granello alla volta, come la sabbia in una
clessidra.

«Parla con lei,» lo esortai, «se devi. Chiedile se è disposta ad
andarsene.»

Sempre che fosse ancora abbastanza vigile da rispondergli. Sempre che si
ricordasse che le avevo parlato.

La sua risposta fu: «La amo più della mia stessa vita.»

Condon chiamò: «Abbiamo bisogno di te quaggiù!»

Bevvi tutto d'un fiato metà della bottiglia mentre Simon mi fissava e le
lacrime gli sgorgavano dagli occhi. L'acqua era pulita e pura, e
deliziosa.

Subito dopo, ero di nuovo con Sorley alle prese con le catene
ostetriche, tirando a tempo con gli spasmi della giovenca gravida e
morente.

ooo

Finalmente, verso mezzanotte estraemmo il vitello, che restò immobile
sulla paglia avviluppato su se stesso, le zampe posteriori ripiegate
sotto il corpo floscio, gli occhi senza vita iniettati di sangue. Condon
rimase in piedi accanto al corpicino per un po'. Poi mi chiese: «C'è
qualcosa che puoi fare?»

«Non posso certo resuscitarlo, se è questo che intendi.»

Sorley mi lanciò uno sguardo d'avvertimento, come per dire: ``Non
torturarlo; tutto questo è già abbastanza difficile''.

Mi avvicinai all'entrata del fienile. Simon era sparito un'ora prima,
mentre stavamo ancora lottando con l'inondazione di sangue emorragico
che aveva impregnato la paglia già fradicia, i nostri abiti, le mani e
le braccia. Dallo spiraglio aperto della porta riuscii a intravedere del
movimento attorno alla macchina - la mia macchina - e un lembo di
tessuto a quadri che poteva essere quello della camicia di Simon.

Stava combinando qualcosa, là fuori. Speravo di sapere cosa.

Sorley spostava lo sguardo dal vitello morto al Pastore Dan Condon,
strofinandosi la barba e senza rendersi conto di impiastricciarla di
sangue. «Forse se lo bruciassimo\ldots» disse.

Condon lo fulminò con un'occhiata raggelante, disperata.

«Ma \emph{forse}» balbettò Sorley.

Poi Simon spalancò le porte del fienile e fece entrare una folata d'aria
fresca. Ci voltammo a guardarlo. La Luna alle sue spalle era gibbosa e
aliena.

«Lei è in macchina» annunciò. «Pronta a partire.» Parlava con me ma
guardava intensamente Sorley e Condon, quasi sfidandoli a replicare.

Il Pastore Dan si limitò a fare spallucce, come se queste faccende
terrene non fossero più rilevanti.

Guardai fratello Aaron. Fratello Aaron si protese verso il fucile.

«Non posso fermarti,» gli dissi, «ma sto per uscire da quella porta.»

Lui si bloccò a metà e aggrottò la fronte. Sembrava stesse cercando di
rimettere a posto la sequenza di eventi che lo avevano portato a quel
momento; ognuno lo aveva condotto al successivo, procedendo con una
sequenza lineare come la salita di una scala, prima un gradino, poi un
altro, poi un altro ancora\ldots{}

Le braccia gli caddero sui fianchi. Si voltò verso il Pastore Dan.

«Penso che se lo bruciamo, comunque, andrà bene lo stesso.»

Mi diressi verso la porta del fienile e raggiunsi Simon, senza voltarmi.
Sorley avrebbe potuto cambiare idea, imbracciare il fucile e prendere la
mira. Non me ne importava davvero più niente.

«Magari lo bruciamo prima dell'alba,» lo sentii dire, «prima che il Sole
spunti di nuovo.»

ooo

«Guida tu» disse Simon quando arrivammo alla macchina. «C'è benzina nel
serbatoio e ce n'è altra dentro alle taniche nel portabagagli. E un po'
di cibo, e altra acqua in bottiglia. Tu guidi e io mi siedo dietro e la
tengo ferma.»

Accesi l'auto e guidai piano sulla salita, oltre la staccionata di
recinzione e l'ocotillo illuminato dalla Luna, verso l'autostrada.
